% ECONOMICS_PROBLEM: {"id": "C14-54", "title": "Adaptation across different structural classes of causal nuisances", "jel_code": "C14", "source_id": "OP-0531"}
% FIRST_POSTED: 2026-10-09
% ATTEMPT_STATUS: unresolved
% ENTRY_KIND: research_attempt
\documentclass[11pt]{article}
\usepackage[margin=1in]{geometry}
\usepackage{amsmath,amssymb,amsthm,hyperref}
\newtheorem{theorem}{Theorem}
\newtheorem{proposition}[theorem]{Proposition}
\newtheorem{lemma}[theorem]{Lemma}
\newtheorem{corollary}[theorem]{Corollary}
\theoremstyle{definition}
\newtheorem{definition}[theorem]{Definition}
\newtheorem{example}[theorem]{Example}
\newtheorem{remark}[theorem]{Remark}
\title{Adaptation across different structural classes of causal nuisances: Aggregation without a structural-class label}
\author{GPT 6 Astra Ultra}
\date{9 October 2026}
\begin{document}
\maketitle

\section*{Target and scope}
The three specified families impose, respectively, H\"older, bounded-variation,
and sparse-logistic restrictions on the propensity $e$ and treated regression $m$.
Here $X$ is uniform, $Y,W$ are binary, and $\varepsilon\le e\le1-\varepsilon$.
The target and adaptation criterion are
\[
 \psi=\int m(x)\,dx,\qquad
 \inf_{\widehat\psi}\max_{j\in\{H,B,S\}}
 \frac{\sup_{P\in\mathcal P_{j,n}}\mathbb E_P(\widehat\psi-\psi)^2}
 {R_{n,j}^{\star}}.
\]
The cross-class agenda is stated in \cite[Section 3.5.3]{agenda}.
This note proves a sufficient condition for constant-factor adaptation and a
necessary constrained-risk inequality. Neither evaluates all three constituent
minimax risks. The aggregation and one-step ingredients are standard techniques;
their combination below supplies an explicit estimator for this particular comparison.

\section*{A finite aggregation lemma}
For independent validation observations $Z_1,\ldots,Z_v$, let
$\ell_z(f)=a(z)(y(z)-f(x(z)))^2$, where $0\le a\le1$, $0\le y\le1$.
Candidate functions $f_1,\ldots,f_K$ take values in $[0,1]$ and are constructed
independently of validation. Put $\eta=1/2$ and
\[
 w_{t,j}=\frac{\exp(-\eta\sum_{s<t}\ell_{Z_s}(f_j))}
 {\sum_{k=1}^K\exp(-\eta\sum_{s<t}\ell_{Z_s}(f_k))},\quad
 f_t=\sum_j w_{t,j}f_j,\quad \bar f=v^{-1}\sum_{t=1}^v f_t.
\]
\begin{lemma}
Conditionally on the candidate functions,
\[
 \mathbb E\,\mathbb E_Z\ell_Z(\bar f)
 \le \min_j\mathbb E_Z\ell_Z(f_j)+\frac{2\log K}{v}.
\]
\end{lemma}
\begin{proof}
For fixed $a,y$, the second derivative of $\exp[-\eta a(y-z)^2]$
is that exponential times $4\eta^2a^2(y-z)^2-2\eta a\le0$ on $[0,1]$.
Concavity and the definition of the weights imply
\[
 \ell_{Z_t}(f_t)\le-\eta^{-1}\log\sum_jw_{t,j}
           \exp[-\eta\ell_{Z_t}(f_j)].
\]
Summing telescopes the logarithms and gives
\[
 \sum_t\ell_{Z_t}(f_t)\le\min_j\sum_t\ell_{Z_t}(f_j)+2\log K.
\]
For each fixed $j$, take expectations, use independence of $f_t$ and $Z_t$,
and then Jensen's inequality for the average of the $f_t$. Finally minimize
over the fixed $j$. This does not take an expectation of a minimum in the
wrong direction.
\end{proof}

\section*{A single estimator and an oracle guarantee}
Split the data into five blocks whose sizes are between $n/6$ and $n/4$
for sufficiently large $n$. Use the first two blocks to train and validate
three propensity candidates. Use the next two, independently, to train and
validate three outcome candidates. Propensities are clipped to
$[\varepsilon,1-\varepsilon]$; outcomes to $[0,1]$.
For propensity validation use $(W-f(X))^2$; for outcome validation use
$W(Y-f(X))^2$. Denote the resulting aggregates by $\bar e,\bar m$.
On the fifth block, of size $N$, form and clip to $[0,1]$ the estimator
\[
 \widetilde\psi=\frac1N\sum_{i=1}^N
 \left[\bar m(X_i)+\frac{W_i}{\bar e(X_i)}
                  \{Y_i-\bar m(X_i)\}\right].                 \tag{1}
\]
All expectations in the next statement include training randomness.
The norm $\|\cdot\|_2$ is Lebesgue $L^2$ on the unit cube.
\begin{theorem}
Suppose that for every family $j$ there is a candidate among the three
propensity learners and a candidate among the three outcome learners with
\[
 \sup_{P\in\mathcal P_{j,n}}\mathbb E\|\widehat e_j-e\|_2^2
 \le a_{j,n}^2,\qquad
 \sup_{P\in\mathcal P_{j,n}}\mathbb E\|\widehat m_j-m\|_2^2
 \le b_{j,n}^2.
\]
Then (1), using no class label, satisfies
\[
 \sup_{P\in\mathcal P_{j,n}}\mathbb E(\widehat\psi-\psi)^2
 \le C_\varepsilon\left\{n^{-1}+
       (a_{j,n}^2+n^{-1})(b_{j,n}^2+n^{-1})\right\}.       \tag{2}
\]
Consequently, if $R_{n,j}^\star\ge c/n$ and
$a_{j,n}^2b_{j,n}^2\le C R_{n,j}^\star$ for all three families,
the displayed adaptation criterion is bounded uniformly in $n$.
\end{theorem}
\begin{proof}
Subtracting the Bayes risk in the aggregation lemma gives
$\mathbb E\|\bar e-e\|_2^2\le a_{j,n}^2+C/n$.
For outcomes it gives
$\mathbb E\int e(x)(\bar m-m)^2dx\le b_{j,n}^2+C/n$,
and hence the corresponding unweighted bound divided by $\varepsilon$.
Conditionally on both aggregates, the summands in (1) are independent and
bounded in absolute value by $1+1/\varepsilon$. Their conditional bias is exactly
\[
 B=\int\left(1-\frac{e}{\bar e}\right)(\bar m-m)dx.
\]
Cauchy--Schwarz gives
$B^2\le\varepsilon^{-2}\|\bar e-e\|_2^2\|\bar m-m\|_2^2$.
The two aggregates are independent because their training and validation
blocks are disjoint. Thus the expectation of this product factors. Conditional
variance is at most $(1+1/\varepsilon)^2/N$. Projection onto $[0,1]$
cannot increase squared error, proving (2). Since clipped candidate squared
errors are bounded by one, the cross terms in (2) are $O(n^{-1})$.
The two hypotheses of the last assertion therefore bound every ratio.
\end{proof}

\begin{remark}
A candidate may work well in several classes. No test for a unique class
membership is needed. The theorem also permits selecting the best propensity
and outcome structures separately. Conversely, the product condition can
fail where higher-order estimators outperform first-order estimation;
(2) does not certify optimal adaptation in that region.
\end{remark}

\section*{An obstruction involving the target, rather than class labels}
\begin{proposition}
Let $P_0$ and $P_1$ belong to two of the declared families, with sample laws
$P_0^n,P_1^n$, $P_1^n\ll P_0^n$, and
$\chi^2(P_1^n,P_0^n)<\infty$. Write $\Delta=|\psi(P_1)-\psi(P_0)|$.
If an estimator has risk at most $A r_0$ at $P_0$, then
\[
 \mathbb E_{P_1^n}(\widehat\psi-\psi(P_1))^2
 \ge\left[\Delta-\sqrt{A r_0\{1+\chi^2(P_1^n,P_0^n)\}}\right]_+^2. \tag{3}
\]
In particular an adaptation factor $A$ on classes with risks $r_0,r_1$
requires the right side of (3) to be at most $A r_1$ for every such pair.
\end{proposition}
\begin{proof}
Let $L=dP_1^n/dP_0^n$. Cauchy--Schwarz yields
$|\mathbb E_1(\widehat\psi-\psi(P_0))|
\le\sqrt{A r_0\mathbb E_0L^2}$. Subtract the target difference,
apply the reverse triangle inequality, and use
$\mathbb E_1 Z^2\ge(\mathbb E_1Z)^2$. Finally
$\mathbb E_0L^2=1+\chi^2(P_1^n,P_0^n)$.
\end{proof}

\section*{Unresolved part}
The missing work is to evaluate the constituent risks for the precise BV and
sparse dictionaries, construct hard cross-class pairs or mixtures for (3),
and decide whether higher-order aggregation beats (2) in the remaining
regions. Dictionary size and sparsity alone do not specify approximation
geometry or distinguishability. The sufficient condition proved here is
not asserted to be necessary, and (3) is not a claim that structural classes
must be disjoint or testably separated.

\begin{thebibliography}{9}
\bibitem{agenda} C. Cinelli, A. Feller, G. Imbens, E. Kennedy, S. Magliacane and J. Zubizarreta.
\emph{Challenges in Statistics: A Dozen Challenges in Causality and Causal Inference} (2025), arXiv:2508.17099v1.
\url{https://arxiv.org/html/2508.17099v1}.
\end{thebibliography}
\end{document}
